\section{SCAN/r2SCAN meta-GGA energy와 gradient}
\label{sec:mgga}

\subsection{Semilocal ingredients}

Spin-unpolarized MGGA는 한 grid point에서
\begin{equation}
  e_{\mathrm{xc}}
  =f(\rho,\sigma,\tau),
  \qquad \sigma=\nabla\rho\cdot\nabla\rho
  \label{eq:mgga-ingredients}
\end{equation}
를 평가한다. Variational derivative에 필요한 weights를
\begin{equation}
  v_\rho=\frac{\partial f}{\partial\rho},
  \qquad
  v_\sigma=\frac{\partial f}{\partial\sigma},
  \qquad
  v_\tau=\frac{\partial f}{\partial\tau}
  \label{eq:mgga-weights}
\end{equation}
로 둔다.

SCAN은 균일 전자 기체, 단일 orbital, 완만히 변하는 density의 limit를 잇도록 설계된
constraint-based functional이다. r2SCAN은 SCAN switching function의 수치적 민감도를
완화해 grid convergence를 개선하지만 $\rho$, $\sigma$, $\tau$를 사용하는 derivative
structure는 같다.

\subsection{AO representation of \texorpdfstring{$\tau$}{tau}}

식~\eqref{eq:tau-definition}을 AO product로 쓰면
\begin{equation}
  \tau(\mathbf r)
  =\frac{1}{2N_k}\sum_{\vk\mu\nu}D_{\mu\nu}^{\vk}
   \sum_{y=x,y,z}
   \partial_y\phi_{\nu\vk}^{*}(\mathbf r)
   \partial_y\phi_{\mu\vk}(\mathbf r).
  \label{eq:mgga-tau-ao}
\end{equation}
Target MGGA Fock의 tau contribution은
\begin{equation}
  V_{\mu\nu}^{\tau,\vk}
  =\frac12\int_\Omega v_\tau(\mathbf r)
   \nabla\phi_{\nu\vk}^{*}(\mathbf r)\cdot
   \nabla\phi_{\mu\vk}(\mathbf r)\dd\mathbf r.
  \label{eq:mgga-vtau-matrix}
\end{equation}
이다. $1/2$는 식~\eqref{eq:tau-definition}에서 온다. Numerical integration library가
이미 $\tau$ convention을 반영한 effective weight를 반환하는지 확인한 뒤 factor를 한
번만 적용한다.

\subsection{Explicit nuclear derivative}

Density matrix를 고정하면
\begin{equation}
  E_{\mathrm{xc}}^{Ax}
  =\int_\Omega\left(
    v_\rho\rho^{Ax}
    +v_\sigma\sigma^{Ax}
    +v_\tau\tau^{Ax}\right)\dd\mathbf r
  \label{eq:mgga-xc-derivative}
\end{equation}
이다. Fixed-cell uniform grid에서는 quadrature point와 volume weight가 핵 좌표에 따라
변하지 않으므로 AO derivative만 남는다. GGA 부분은
\begin{equation}
  \sigma^{Ax}=2\nabla\rho\cdot\nabla\rho^{Ax}
  \label{eq:mgga-sigma-derivative}
\end{equation}
이고 tau derivative는
\begin{align}
  \tau^{Ax}(\mathbf r)
  =\frac{1}{2N_k}\sum_{\vk\mu\nu}D_{\mu\nu}^{\vk}
   \sum_y\big[&
   \partial_y\phi_{\nu\vk}^{Ax *}
   \partial_y\phi_{\mu\vk}
  \notag\\
   &+\partial_y\phi_{\nu\vk}^{*}
   \partial_y\phi_{\mu\vk}^{Ax}\big].
  \label{eq:mgga-tau-derivative}
\end{align}
AO center derivative와 spatial gradient가 함께 있으므로 second spatial derivative AO가
필요하다. 이 점이 GGA보다 MGGA gradient kernel이 복잡한 주된 이유다.

\subsection{Nonvariational response와의 결합}

MGGA에서도 source orbital은 HF가 결정한다. 따라서 Hessian $A$는 식
\eqref{eq:hf-hessian-action}과 동일하고, target Fock에 식
\eqref{eq:mgga-vtau-matrix}가 더해져 target VO gradient
\begin{equation}
  g_{ai}^{\vk,\mathrm{MGGA}}
  =\left[C_v^{\vk\dagger}
   \left(h+J+V_\rho+V_\sigma+V_\tau\right)^{\vk}
   C_o^{\vk}\right]_{ai}
  \label{eq:mgga-target-vo}
\end{equation}
를 바꾼다. Z-vector solve 후 식~\eqref{eq:response-gradient}에 source HF nuclear
perturbation을 contraction한다. MGGA kernel을 explicit XC term에만 연결하고 target Fock
VO RHS를 PBE 값으로 남기면 nonvariational response가 틀린다.

\subsection{Grid와 aliasing}

Gaussian AO의 high-frequency component는 grid가 거칠면 aliasing된다. $\tau$는 AO
gradient product이므로 density보다 높은 reciprocal component에 민감하다. 따라서
MGGA에서 \texttt{cell.mesh} 또는 effective integration mesh convergence를 별도로 확인해야
한다. Mesh를 한 방향씩 바꾸지 말고 resolved integer triplet과 cutoff를 함께 기록한다.

Translation sum rule은 grid aliasing을 민감하게 보여준다. Analytic--FD component error가
작아도
\begin{equation}
  \left\|\sum_A\mathbf g_A^{\mathrm{xc}}\right\|
  \label{eq:mgga-xc-translation}
\end{equation}
가 크면 explicit XC grid derivative가 아직 수렴하지 않은 것이다.

\subsection{CPU/GPU parity}

CPU와 GPU는 같은 libxc functional label, density convention, mesh, AO derivative order를
사용해야 한다. 비교 항목은 total energy와 final gradient뿐 아니라
$E_{\mathrm{xc}}$, target Fock VO block, $v_\tau$ row derivative, Z-vector RHS와 residual이다.
SCAN/r2SCAN completion case에는 Gamma H$_2$, distorted H$_2$O와 sampled k-point case가
포함되어야 한다. 또한 target MGGA와 source-HF density-fitting provider는 서로 다른
책임을 가지므로 FFTDF와 GDF만이 아니라 RSDF에서도 같은 $\rho$, $\nabla\rho$, $\tau$
target path를 검증해야 한다. RSDF+MGGA case는 Eq.~\eqref{eq:mgga-xc-derivative}의
XC derivative와 Eq.~\eqref{eq:rsdf-total-gradient}의 source/metric auxiliary response가
동시에 작동하는 composition test이다.
